\chapter{Version 1: Random Forest}\label{ch:v1}

\section{Data Preparation}

The Version 1 data pipeline uses StatsBombPy~\cite{statsbombpy} to extract events from all available matches in the open dataset. Events are filtered to passes, carries, and shots, and their scalar features are extracted into a flat CSV file (\texttt{statsbomb\_chained\_dataset.csv}). Possession chains are identified and labelled as described in Section~\ref{sec:chains}, with goal-chain events receiving a label of 1 and all others receiving 0.

The dataset is split at the match level using a 70/15/15 train/validation/test partition with a fixed seed (42). Version 1 does not use a validation loop during training (Random Forests have no early stopping by design~\cite{breiman2001}), so the validation matches are withheld from training but not actively used during it.

Team form is computed using an exponentially decayed rolling window over the team's last five matches, giving more weight to recent results. Specifically, with decay factor $\lambda = 0.75$, the most recent match receives weight 1.0, the previous match 0.75, and the oldest match in the five-game window approximately 0.32 ($\lambda^4 = 0.75^4 \approx 0.316$). The full formulation is given in Chapter~\ref{ch:features}. This contrasts with a simple unweighted average and reflects the intuition that recent form is a stronger predictor of current performance than older results.

\section{Model Architecture and Design Choices}

Version 1 uses a \texttt{RandomForestClassifier} with 100 trees and a maximum depth of 12 (see Chapter~\ref{ch:methodology} for how Random Forests work). A Random Forest was chosen as the Version 1 baseline for three reasons: (i) it yields interpretable feature importances directly from the tree structure, which inform the feature selection choices in Versions 2 and 3 (Section~\ref{sec:v1_importance}); (ii) it is robust to differences in feature scale without requiring normalisation; and (iii) it is a well-established tabular baseline in football analytics~\cite{lucey2015}. A gradient-boosted ensemble such as XGBoost would likely produce a marginally higher AUC, but at the cost of interpretability and without the feature-importance output needed to motivate the Version 2 architecture.

The maximum depth of 12 is a regularisation parameter: unconstrained trees would fit the training data perfectly but overfit badly on unseen events. Depths from 6 to unconstrained were evaluated on the validation set; depth 12 produced the highest validation AUC without signs of overfitting.

Class imbalance is not explicitly addressed in Version 1. The Random Forest uses equal sample weights, which means the model can partially compensate by learning a conservative low-probability prior for most events.

\section{Results}\label{sec:v1_results}

Version 1 achieved a test ROC AUC of 0.6840, Brier Score of 0.0161, and Log Loss of 0.0802. The ROC AUC of 0.68 is a moderate improvement over random (0.5), confirming the model can distinguish goal-chain events from background possession. 

The heatmap (Figure~\ref{fig:v1_heatmap}) confirms that the model has learned a broadly sensible spatial representation: threat is concentrated in and around the attacking penalty area, peaks just in front of the six-yard box, and drops off sharply as the ball moves deeper or wider. The defensive half is effectively zero throughout.

The decision-tree structure of the model is also visible: rather than a smooth radial gradient, the threat surface has jagged, stepped edges and a handful of isolated blobs around the midfield area. These are artefacts of the axis-aligned decision boundaries that an ensemble of trees produces ; a CNN, which learns spatial filters directly from data, does not have this limitation.

\section{Feature Importance}\label{sec:v1_importance}

The feature importance results (Figure~\ref{fig:v1_importance}) reveal a clear hierarchy. \texttt{dist\_to\_goal} is overwhelmingly dominant at 0.393, accounting for nearly 40\% of total impurity reduction. \texttt{start\_x} is a distant second at 0.157, followed by \texttt{start\_y} and \texttt{angle\_to\_goal} which are essentially tied at 0.096 each.

The three main context features contribute meaningfully: \texttt{minute} scores 0.080, \texttt{score\_diff} 0.068, and \texttt{team\_form} 0.056 , together accounting for roughly 20\% of total importance. This confirms that game-state context is a genuine signal, even if secondary to position. The action-type flags (\texttt{type\_Pass}, \texttt{type\_Carry}) contribute the least, as position already encodes most of the relevant information once action type is conditioned on.

\section{Limitations and Motivation for Version 2}

Distance to goal alone accounts for nearly 40\% of the model's predictive power, and the top four spatial features together account for around 74\%. The model is essentially learning a distance-to-goal heuristic with a contextual adjustment from match state. More fundamentally, it has no access to the positions of any of the other 21 players on the pitch. A flat scalar feature vector cannot encode the spatial arrangement of defenders and attackers, and a Random Forest cannot learn the complex geometric relationships between them from a list of numbers. Incorporating full player positional context requires a richer input representation.
